\section*{Part II: Choosing and Governing the Legal Structure}
\addcontentsline{toc}{section}{Part II: Choosing and Governing the Legal Structure}

\subsection{7. A private limited company: the conventional commercial vehicle}

An Indian private company limited by shares is often suitable for a scalable commercial education, software or consulting enterprise with employees, investors and share-based ownership. Its benefits may include limited liability, transferable economic interests subject to rules, defined director and shareholder governance, and compatibility with many financing instruments. It brings statutory filings, accounts, board duties, tax and employment obligations.

Limited liability does not protect a founder from every personal guarantee, statutory breach, fraud, wrongful conduct or independently assumed contractual liability. Treat the corporate boundary as real in accounting and conduct.

\subsection{8. LLP, sole enterprise and other forms}

An LLP may suit some professional services and internal profit-sharing arrangements but can be less convenient for venture-style equity and employee options; exact rules differ. A sole proprietorship offers simplicity but weak separation of business obligations. Partnerships and trusts involve different liability and governance structures. A founder should choose based on customers, risk, hiring, investor needs, taxation and future statutory activities - not the lowest incorporation fee.

Never convert legal entities or transfer IP casually; conversions can trigger tax, stamp, consent and contract issues.

\subsection{9. Section 8 and other mission structures}

An Indian Section 8 company can be established for qualifying public-interest purposes, including science, education and research, under the Companies Act and relevant rules. It is not a route for distributing profits as ordinary shareholder dividends; the legal restrictions and governance must be understood. Trusts and societies are other possible mission vehicles, with state-specific laws and sponsor eligibility questions. MCA's incorporation materials describe Section 8 objects and procedural particulars [G1].

A philanthropic research center, university sponsor and commercial company may need different structures. Do not assume 1 entity can freely mix VC profits, gifts, degree fees, research grants and a founder-owned property portfolio.

\subsection{10. Entity architecture follows obligations, not branding}

A founder may eventually operate a commercial company, separately governed educational/research entity, campus property holder, and optional venture spinouts. Multiple entities are justified when their legal obligations, risk pools, beneficiaries, asset rights and capital sources genuinely differ. They are not justified merely to multiply prestigious letterheads or move funds opaquely.

Draw a diagram with every entity, its owners/members, assets, liabilities, decision makers, intercompany agreements, tax registrations and restrictions on funds. If you cannot explain money and rights moving between them, the structure is not ready.

\subsection{11. The company register is the source of corporate truth}

Maintain incorporation documents, memorandum and articles, registers of members and directors, share allotment records, board and shareholder resolutions, beneficial ownership records, loan agreements, statutory filings, material contracts, intellectual-property assignments and employment agreements. Use a qualified company secretary where appropriate. A cap table in a slide deck is not a substitute for legal share records.

Before raising funds, reconcile the \textbf{issued shares}, paid-up capital, securities outstanding, vesting and option plans, convertible rights, and actual statutory filings. Investors may discover inconsistencies that undermine valuation or closing.

\subsection{12. Founders' agreements before disputes}

Co-founders should clarify ownership, cash and time contributions, vesting or reverse-vesting structures where legally and contractually feasible, roles, compensation, IP assignment, confidentiality, departure and incapacity, transfer restrictions, deadlock resolution, decision rights, dispute resolution and what happens when someone ceases contributing. The agreement must align with constitutional documents and applicable law.

A handshake 50/50 split can become an unmanageable veto problem if 1 founder leaves. Equally, ``I started the idea'' is not a legal substitute for properly executed share and IP arrangements.

\subsection{13. Intellectual property ownership is a governance decision}

For software, curriculum, research, trademarks, datasets and inventions, record the author, contracting relationship, assignment or license, permitted uses and 3rd-party restrictions. Former employer IP, open-source dependencies and contractor-created code can be expensive diligence problems. A developer who receives a fee does not necessarily convey every intellectual-property right absent suitable legal agreements.

Use rights-cleared materials and explicit independent work boundaries. For research partnerships, establish background/foreground IP and publication rights before results exist.
